\section{The pure-state factor-two theorem}\label{sec:factor}
The two-sided mixed-state comparison and the pure-density comparison use
different pairs of spectral measures. The latter has a tensor structure
that protects a universal inequality.

\begin{theorem}[Finite-dimensional factor-two inequality]\label{thm:factor}
Let $G$ be a Hermitian matrix and $\psi$ a unit vector in a finite-dimensional
Hilbert space. Set $\psi(t)=e^{-itG}\psi$ and
$\sigma(t)=\ket{\psi(t)}\bra{\psi(t)}$. With the state Krylov chain
generated by $(G,\psi)$ and the operator chain generated by
$([G,\cdot],\sigma(0))$, for every real $t$,
\begin{equation}\label{eq:factor-two}
 \CK(\sigma(t))\ge2\CS(\psi(t)).
\end{equation}
The coefficient two is sharp. Stationary seeds, repeated eigenvalues, and
terminating chains are included.
\end{theorem}
\begin{proof}
Under row vectorization,
\begin{equation}\label{eq:factor-vectorization}
 \chi=\psi\otimes\bar\psi,\qquad
 L=G\otimes I-I\otimes\bar G,\qquad
 \chi(t)=\psi(t)\otimes\overline{\psi(t)}.
\end{equation}
Let $K_i$ be the state Krylov basis and define the polynomial subspaces
\begin{equation}\label{eq:flags}
 V_n=\spanof\{L^k\chi:0\le k\le n\},\qquad
 W_n=\spanof\{K_i\otimes\bar K_j:i+j\le n\}.
\end{equation}
The binomial expansion of $L^k$ proves $V_n\subseteq W_n$. This inclusion
remains valid if a chain terminates, or if the joint cyclic subspace is a
proper subspace of the tensor-product space. No completion of the joint
Krylov basis is needed.

Write $p_i=|\bra{K_i}\psi(t)\rangle|^2$ for the state probabilities and
$p_k^{\rm op}$ for the operator probabilities. Conjugate evolution has the
same probabilities; the possible signs $(-1)^j$ of its Krylov vectors do
not affect their spans. Nested orthogonal projections give
\begin{equation}\label{eq:tail-domination}
 \sum_{k>n}p_k^{\rm op}
 =1-\norm{P_{V_n}\chi(t)}^2
 \ge1-\norm{P_{W_n}\chi(t)}^2
 =\sum_{i+j>n}p_i p_j.
\end{equation}
Summing these nonnegative tails yields
\begin{equation}\label{eq:tail-sum}
 \CK(\sigma(t))\ge\sum_{i,j}(i+j)p_i p_j
 =2\sum_i i p_i=2\CS(\psi(t)).
\end{equation}
For any nonstationary finite-dimensional seed, with
$v=\bra\psi(G-\langle G\rangle)^2\ket\psi>0$, the short-time
coefficients of $\CK$ and $\CS$ are $2v$ and $v$. Their ratio tends to two,
so no larger universal multiplier is possible. A stationary seed gives
zero on both sides.
\end{proof}

The proof also identifies the excess complexity exactly:
\begin{equation}\label{eq:flag-gap}
 \CK(\sigma(t))-2\CS(\psi(t))
 =\sum_{n\ge0}\norm{(P_{W_n}-P_{V_n})\chi(t)}^2.
\end{equation}
In finite dimension the sum terminates. Its positivity is an assertion
about the evolved cyclic vector and the nested subspaces. It does not
require an unqualified positivity claim for a joint number operator on an
unused ambient complement.

\begin{corollary}[The two purification branches]\label{cor:branches}
For either of the source's unitary purification schemes on a fixed finite
enlarged Hilbert space, using its actual normalized initial purification,
\begin{equation}\label{eq:branches}
 \KI(t)\ge2\SI(t),\qquad \KU(t)\ge2\SU(t)
 \qquad(t\in\mathbb R).
\end{equation}
\end{corollary}
\begin{proof}
Apply Theorem~\ref{thm:factor} first with $G=G_I$, and separately with
$G=G_{U^*}$. Both use the pure seed $s$ and the rank-one operator seed
$ss^\dagger$. The two generators need not have the same Krylov structure.
The theorem itself does not require full rank of the original density matrix.
\end{proof}

Equation~\eqref{eq:factor-two} is a special case of the prior
tensor-product superadditivity theorem of Murugan and van Zyl
\cite{MuruganVanZyl2026}: use generators $G$ and $-\bar G$ and the product
seed $\psi\otimes\bar\psi$. Their two subsystem complexities are equal.
The self-contained subspace proof above fixes the source matching without
making the factor-two theorem a novelty claim. We restrict the theorem here
to finite dimension; no infinite-dimensional extension is required for any
counterexample or no-go result in this paper.
